\SourceLeaf{097}
\[
u=22877^{\mathrm m}94\times\frac{\sin(v^{(13)}-v^{(12)}-0''652).\sin(v^{(14)}-v^{(16)}-0''814)}{\sin(v^{(1)}-v^{(0)}-0''652).\sin(v^{(6)}-v^{(4)}-0''814)}
\]

Huius valor e valoribus correctis directionum \(v^{(0)},v^{(1)}\) etc. invenitur
\[
=26766^{\mathrm m}68
\]

Differentiatio autem illius expressionis suppeditat, si differentialia \(\mathrm dv^{(0)},\allowbreak\mathrm dv^{(1)}\) etc. minutis secundis expressa concipiuntur,
\[
\begin{aligned}
\mathrm du={}&0^{\mathrm m}16991(\mathrm dv^{(0)}-\mathrm dv^{(1)})+0^{\mathrm m}08836(\mathrm dv^{(4)}-\mathrm dv^{(6)})\\
&-0^{\mathrm m}03899(\mathrm dv^{(12)}-\mathrm dv^{(13)})+0^{\mathrm m}16731(\mathrm dv^{(14)}-\mathrm dv^{(16)})
\end{aligned}
\]

Hinc porro invenitur
\[
\begin{aligned}\relax
[al]&=-0{,}08836\\
[bl]&=+0{,}13092\\
[cl]&=-0{,}00260\\
[dl]&=+0{,}07895\\
[el]&=+0{,}03899\\
[fl]&=-40{,}1315\\
[gl]&=+10{,}9957\\
[ll]&=+0{,}13238
\end{aligned}
\]

Hinc denique per methodos supra traditas invenitur, quatenus metrum pro unitate dimensionum linearium accipimus,
\[
\frac1P=0{,}08329,\quad\text{sive}\quad P=12{,}006
\]
unde error medius in valore lateris Falkenberg-Breithorn metuendus \(=0{,}2886m\) metris, (ubi \(m\) error medius in directionibus observatis metuendus, et quidem in minutis secundis expressus), adeoque, si valorem ipsius \(m\) supra erutum adoptamus,
\[
=0^{\mathrm m}1209
\]

Ceterum inspectio systematis triangulorum sponte docet, punctum Hauselberg omnino ex illo elidi potuisse, incolumi manente nexu inter latera Wilsede-Wulfsode atque Falkenberg-Breithorn. Sed a bona methodo abhorret, \textit{supprimere} idcirco observationes, quae ad punctum Hauselberg referun\SourceLeaf{098}tur\footnote{Maior pars harum observationum iam facta erat, antequam punctum Breithorn repertum, atque in systema receptum esset.}, quum certe ad praecisionem augendam conferre valeant. Ut clarius appareret, quantum praecisionis augmentum inde redundet, calculum denuo fecimus excludendo omnia, quae ad punctum Hauselberg referuntur, quo pacto e 18 directionibus supra traditis octo excidunt, atque reliquarum errores maxime plausibiles ita inveniuntur:
\[
\begin{array}{r@{\;}l@{\qquad}|@{\qquad}r@{\;}l}
(0)=&+0''327&(12)=&+0''206\\
(1)=&-0{,}206&(13)=&-0{,}206\\
(3)=&-0{,}121&(14)=&+0{,}327\\
(4)=&+0{,}121&(15)=&+0{,}206\\
(6)=&-0{,}121&(16)=&+0{,}121
\end{array}
\]
Valor lateris Falkenberg-Breithorn tunc prodit \(=26766^{\mathrm m}63\), parum quidem a valore supra eruto discrepans, sed calculus ponderis producit
\[
\frac1P=0{,}13082\quad\text{sive}\quad P=7{,}644
\]
adeoque error medius metuendus \(=0{,}36169m\) metris \(=0^{\mathrm m}1515\). Patet itaque, per accessionem observationum; quae ad punctum Hauselberg referuntur, pondus determinationis lateris Falkenberg-Breithorn auctum esse in ratione numeri 7,644 ad 12,006, sive unitatis ad 1.571.
\par\medskip\begin{center}\rule{.28\textwidth}{.3pt}\end{center}
